\section{Retour observable et persistance du registre}
\label{nuclear:9}
\subsection{Récupération conjointe à partir des retours complets}

Soit \(X\) le domaine des histoires complètes étalonnées de
\cref{ret:extension}, y compris l’origine du calendrier et, lorsqu’il est
utilisé, l’indice discret du système de coordonnées. Les états conservent les
transitions et le registre de mémoire. Nous écrivons
\(r=\operatorname{Ret}_+:X\to Y_+\) pour l’échantillonnage aux retours
complets, avec \(Y_+=r(X)\).

Le \cref{ret:homeomorfismo} fournit son inverse explicite :
\(x_{3j+s}=F^s(y_j)\), pour \(s=0,1,2\). La reconstruction utilise
l’état complet conservé à chaque retour et la mise à jour
effective ; \eqref{ret:conjugacion} conserve également le transport.

Soient \(a:X\to R_{\rm HMT}\) l’évaluation numérique par cylindres et
\(q:X\to E_\infty\) l’incidence avec sa jauge, selon
\cref{sec:generacion,exc:doble-lectura,exc:limite-inverso}.
Soit \(d:X\to X\) un transport réversible et continu, avec la jauge
transportée. Sur \(Y_+\), on définit




\[
a_+=a\circ r^{-1},\qquad q_+=q\circ r^{-1},\qquad d_+=r\circ d\circ r^{-1}.
\]



Alors



\begin{equation}
\boxed{(a,q)=(a_+,q_+)\circ r,\qquad rd=d_+r.}
\label{nuc09:lecturas}
\end{equation}



\begin{proof}
C’est le \cref{ret:lecturas} : les compositions avec \(r\)
et \(r^{-1}\) s’annulent, et l’inverse de \(d_+\) est \(rd^{-1}r^{-1}\).
Les deux lecteurs sont transportés à partir du même état complet.
\end{proof}

En particulier, pour \(N\in\ZZ_{>0}\),
\(d^Nx=x\) si et seulement si \(d_+^Nr(x)=r(x)\).
La conjugaison conserve la périodicité ou l’apériodicité de l’état
complet. La dynamique d’un lecteur partiel est déterminée au moyen
de son propre facteur de transport.

\subsection{Représentation ultérieure explicite du transport
}

Sur l’ensemble des histoires déjà engendré, on considère la réalisation
ultérieure \(H=\ell^2(X)\), de base orthonormale
\((\delta_x)_{x\in X}\) et munie de la mesure de comptage. Chaque vecteur a un support
au plus dénombrable. Le transport réversible détermine



\[
C\delta_x=\delta_{d x},\qquad C^{-1}\delta_x=\delta_{d^{-1}x}.
\]



La bijection \(d\) prouve que \(C\) est unitaire.
Cette représentation conserve les étiquettes complètes des histoires.
L’espace \(X\) et son transport précèdent cette réalisation linéaire.


Le système de coordonnées nonadique de \cref{nuclear:2} détermine



\[
T=\frac{8I+C}{9},\qquad
\eta v=\frac1{27}(8(C-I)v,-(C-I)v,\ldots,-(C-I)v),
\]





\[
I-T^*T=\eta^*\eta=\frac8{81}(I-C)^*(I-C).
\]



Soit \(P\) la projection orthogonale sur \(\ker(I-C)\).
La preuve spectrale de \cref{nuclear:2,nuclear:3} donne
\(T^N\to P\) fortement. Le télescopage fournit l’isométrie



\[
\mathcal Vv=(Pv,\eta v,\eta Tv,\eta T^2v,\ldots),\qquad
\mathcal M=\operatorname{Im}\mathcal V.
\]



Le décalage de registres \(B\) conserve la première composante
et élimine le premier bloc de mémoire. Sur l’image compatible,




\begin{equation}
\boxed{\widehat C=\mathcal VC\mathcal V^{-1}
=(9B-8I)|_{\mathcal M}.}
\label{nuc09:transporte-recuperado}
\end{equation}



L’inverse de la représentation est donné par



\begin{equation}
v=Pv+\sum_{j\ge0}T^{*j}\eta^*(\eta T^jv).
\label{nuc09:inversa-registro}
\end{equation}



\begin{proof}
L’identité de défaut se télescope jusqu’à \(N\) ; sa limite conserve
exactement la norme. Les identités
\(\mathcal VT=B\mathcal V\) et \(C=9T-8I\) donnent
\eqref{nuc09:transporte-recuperado}.
L’isométrie \(\mathcal V\) est surjective sur \(\mathcal M\) ;
prendre son adjoint donne \eqref{nuc09:inversa-registro}, avec convergence
en norme de la série. L’unitarité du transport retrouvé
est établie dans \(\mathcal M\).
\end{proof}

Pour \(A\in\mathcal B(H)\), soit
\(\widehat A=\mathcal VA\mathcal V^{-1}\) dans \(\mathcal M\).
Cette correspondance préserve la norme, les produits, les adjoints et les commutateurs,
car \(\mathcal V^{-1}=\mathcal V^*|_{\mathcal M}\).
Elle représente conjointement les observables numériques, ceux d’incidence
et leurs relations avec \(C\).

\subsection{Convergence numérique et corrélations avec l’incidence
}

Les cylindres de \eqref{eq:cilindro} fournissent, dans le système de coordonnées
d’indice entier \(m(x)\), six trits par bloc. Nous écrivons



\[
\xi_n(x)=\sum_{k=1}^{6n}d_k(x)3^{-k},\qquad
\xi(x)=a(x)-m(x),\qquad 0\le\xi(x)-\xi_n(x)\le3^{-6n}.
\]



Ici, \(d_k\in\{0,1,2\}\) est la lecture ternaire ultérieure du trit ;
son orientation équilibrée conserve son typage propre. Le système de coordonnées
retient \(m\) séparément, y compris la convention des extrémités,
et reconstruit \(a=m+\xi\). Aux extrémités de double représentation,
\(\xi\) est interprété dans ce système de coordonnées.

Soit \(M_f\) la multiplication par \(f\) dans \(\ell^2(X)\).
La normalisation \(0\le\xi_n\le\xi\le1\) fait que
\(M_\xi\) et \(M_{\xi_n}\) sont des contractions.
Pour un ensemble \(B\) du codomaine d’incidence fini de \(q_j\),
nous définissons
\(E_{j,B}=M_{\mathbf1_{q_j^{-1}(B)}}\).
C’est le projecteur sur les histoires qui satisfont cette incidence.

\begin{theorem}[Conservation simultanée de la précision et du transport]
\label{nuc09:precision}
Pour tout \(t\in\ZZ\),



\begin{equation}
\boxed{
\left\|\widehat C^{-t}\widehat M_\xi\widehat C^t
-\widehat C^{-t}\widehat M_{\xi_n}\widehat C^t\right\|
=\|M_\xi-M_{\xi_n}\|\le729^{-n}.}
\label{nuc09:error-observable}
\end{equation}



Les relations entre les projecteurs incidenciels, leurs produits
et leur transport sont conservées exactement.
Pour un produit ordonné de contractions contenant \(r\)
observables numériques \(M_\xi\), évalués à des temps entiers
arbitraires, et un nombre fini quelconque de facteurs incidenciels
ou unitaires exacts, remplacer ces \(r\) facteurs par
\(M_{\xi_n}\) change le produit en norme d’au plus



\begin{equation}
\boxed{r\,729^{-n}.}
\label{nuc09:error-producto}
\end{equation}



La même borne vaut pour tout élément de matrice entre vecteurs
unitaires. Les facteurs peuvent être non commutatifs ; leur ordre est conservé.

\end{theorem}

\begin{proof}
La norme d’un multiplicateur sur \(\ell^2(X)\) est
\(\sup_{x\in X}|f(x)|\) ; la queue ternaire prouve la première borne.
Conjuguer par \(C^t\) et par \(\mathcal V\) conserve la norme.
Pour les produits \(A_1\cdots A_s\) et \(B_1\cdots B_s\), on utilise



\[
\prod_{i=1}^s A_i-\prod_{i=1}^s B_i
=\sum_{j=1}^s B_1\cdots B_{j-1}(A_j-B_j)A_{j+1}\cdots A_s.
\]



Les facteurs exacts donnent une différence nulle ; chacun des \(r\)
restants apporte une norme au plus égale à \(729^{-n}\).
La sous-multiplicativité et la borne un des autres facteurs prouvent
\eqref{nuc09:error-producto}.
Cauchy--Schwarz donne l’affirmation sur les éléments de matrice.

\end{proof}

Cette uniformité conserve la précision lors du transport de l’observable
approché complet. L’opérateur \(C\) et l’incidence restent
exacts dans \eqref{nuc09:error-observable}.
L’indice \(n\) mesure la profondeur de cylindre ; \(t\) mesure les itérations
du transport \(d\) ; \(j\), dans \eqref{nuc09:inversa-registro},
mesure les itérations de compression.

\subsection{Covariance de l’opérateur de mémoire sous le retour}

On restreint \(X\) et sa représentation atomique au secteur réversible
gradué \(X_{\rm rev,grad}\) de l’horloge de
\eqref{eq:weyl-relojes}. On conserve dans chaque histoire le degré entier
de mémoire \(q_{\rm mem}\). Pour \(d=\Gamma_9\), sa mise à jour est
\(q_{\rm mem}(dx)=q_{\rm mem}(x)+1\).
L’opérateur de mémoire \(D\) est la multiplication par
\(q_{\rm mem}\), de domaine



\[
\operatorname{Dom}D=\{v:\sum_{x\in X}|q_{\rm mem}(x)|^2|v_x|^2<\infty\}.
\]



Le changement d’indice \(x\mapsto dx\) et les inégalités
\((q\pm1)^2\le2q^2+2\) prouvent que \(C\) et \(C^{-1}\)
conservent ce domaine. Sur celui-ci,




\begin{equation}
C^*DC=D+I,\qquad
\widehat C^*\widehat D\widehat C=\widehat D+I_{\mathcal M},
\qquad\operatorname{Dom}\widehat D=\mathcal V\operatorname{Dom}D.
\label{nuc09:memoria}
\end{equation}



\begin{proof}
Sur \(\delta_x\), le premier membre vaut
\(q_{\rm mem}(dx)\delta_x=(q_{\rm mem}(x)+1)\delta_x\).
L’identité s’étend au domaine indiqué par troncature du
support dans la norme du graphe de \(D\).
La seconde égalité s’obtient par conjugaison unitaire vers
\(\mathcal M\).
\end{proof}

Par récurrence,
\(q_{\rm mem}(d^Nx)=q_{\rm mem}(x)+N\).
Chaque entier \(N>0\) change l’état complet.
Toute orbite de \(d\) est infinie ; une fonction de carré sommable
constante sur chacune de ces orbites est nulle.
Dans cette réalisation bilatérale, \(P=0\).
La formule \eqref{nuc09:inversa-registro} reconstruit alors
l’état linéaire complet à partir de ses compléments de mémoire.
Cette conclusion utilise la graduation effective de l’horloge.


Le retour de phase coexiste avec l’incrément de mémoire de
\eqref{nuc09:memoria}.
L’apériodicité sturmienne et sa complexité appartiennent au facteur
symbolique de \cref{sec:generacion} ; le présent résultat conserve
la dynamique graduée de l’état complet.

\subsection{Incidence, raffinement et trois régimes}

Le registre \(K\) possède en outre la réalisation complète de
\cref{nuclear:5},
\(\mathcal Fk=(\mu(k),IP_0k/6)\).
Son inverse conserve les douze coordonnées, y compris la moyenne.
Appliquer \(\mathcal F\) aux fibres correspondantes et conserver
les étiquettes d’histoire transporte cette représentation.
L’application \(\mathcal F\circ K\) conserve exactement les fibres
du lecteur \(K\) ; la récupération des histoires complètes provient
de \(r\) ou de ses étiquettes effectives.

Les systèmes de raffinement isométrique de
\cref{nuclear:1,nuclear:3,nuclear:5} conservent
\(\mathcal V\) et \(C\) au moyen de leurs carrés d’entrelacement.
La preuve utilise les réalisations par niveau qui y sont construites.
La représentation atomique \(\ell^2(X)\) et les réalisations
\(L^2\) à mesure projective conservent leurs domaines
et mesures correspondants.

Le lien avec le pliage et le dépliage est développé dans
\cref{euler-trit-generadores,euler-trit-evaluaciones} :
Coxeter, transport nilpotent et action de \(F_{\rm av}^2\).
Pour chaque réalisation bidimensionnelle \(E\) qui conserve la forme
alternée \(\Omega\), le système de coordonnées nonadique satisfait



\begin{equation}
\Omega=T_E^*\Omega T_E+D_E^*\Omega D_E,
\qquad T_E=(8I+E)/9,\quad D_E=\sqrt8(E-I)/9.
\label{nuc09:balance-alternante}
\end{equation}



Le développement algébrique annule les termes croisés et utilise
\(E^*\Omega E=\Omega\). Dans la réalisation réelle,
\(*\) signifie transposée.
La loi orientée réunit les trois régimes au moyen des formes
et des transports de \cref{nuclear:7,nuclear:8}.
Sa réalisation comme action utilise en outre les normalisations
et les connexions qui y sont spécifiées.

La relation avec la dynamique de connexion se formule au moyen
des constructions de \cref{nuclear:6,nuclear:7,nuclear:8}.
Les identités
\eqref{nuc09:transporte-recuperado}--\eqref{nuc09:balance-alternante}
conservent leur représentation, la précision des observables
et les formes transportées.

